% Volume VII: Formal Synthesis
% Part XIV. Dynamics, Geometry, and Holonomy
% Chapter 79: Curvature of Reinterpretation
% Spine: proof-spines/ch079-spine.md

\chapter{Curvature of Reinterpretation}
\label{ch:ch079}

This chapter makes order-sensitive judgment explicit through the \textbf{curvature of reinterpretation}. For update operators \(\nabla_i\) and \(\nabla_j\), curvature is registered by the commutator
\[
[\nabla_i,\nabla_j]J
=
\nabla_i\nabla_j J-\nabla_j\nabla_i J.
\]
\Definition\ A nonzero value means that applying two evidentiary or interpretive updates in opposite orders yields different resulting judgment \(J\). Sequence, framing, and procedural posture are therefore not external accidents; they are formally encoded in the geometry of revision.

The chapter's claim is modest but sharp. When the commutator does not vanish, reinterpretation is locally curved rather than flat: nearby judgments do not translate rigidly under new evidence, but bend differently depending on the route by which the evidence arrives. Legal and institutional order-sensitivity can thus be represented geometrically, setting up the field-theoretic account of public interpretation that follows.
